\begin{theorem}[derived global CDP criterion]
\label{thm:corrected-global}
Let $X$ be a connected compact complex threefold with
\[
 H^1(X;\Z)=H^2(X;\Z)=0,\qquad a(X)=1,
 \tag{5.3}
\]
and suppose its algebraic reduction is a holomorphic map $f:X\to C$ to a
smooth compact curve.  If there is a line bundle $L_0\in\Pic(X)$ with
\[
 H^2(X,T_X\otimes L_0)=0,
 \tag{5.4}
\]
then a dense open subset $U\subset\Pic(X)$ satisfies
\[
  H^i(X,T_X\otimes M)=0\quad(i\in\{0,2,3\},\ M\in U),
 \qquad
 \left\langle c_3(T_X),[X]\right\rangle\le0.
 \tag{5.5}
\]
\end{theorem}

\begin{proof}
The exponential sequence and (5.3) identify
\(\Pic(X)\) with the complex vector group $H^1(X,\cO_X)$; in particular it
is connected and irreducible and is stable under dualization.  Choose a Stein
cover $\{U_i\}$ of $X$ and a \v{C}ech basis
$[a^{(1)}_{ij}],\ldots,[a^{(q)}_{ij}]$ of $H^1(X,\cO_X)$ represented by additive
cocycles.  If $\lambda=(\lambda_1,\ldots,\lambda_q)$ are the corresponding
coordinates, the functions
$\exp(\sum_r\lambda_r a^{(r)}_{ij})$ satisfy the multiplicative cocycle
condition and hence define the holomorphic Poincar\'e family on
$X\times\Pic(X)\simeq X\times H^1(X,\cO_X)$.  Every
line bundle in it has vanishing integral
first Chern class.  The algebraic-reduction map is surjective; choose
$c\in C$ and let $D=f^*(c)$ denote the scheme-theoretic fibre divisor.  It is
a nonzero effective Cartier divisor and
\(\cO_X(D)=f^*\cO_C(c)\in\Pic^0(X)\), because $H^2(X;\Z)=0$.
  The stabilization argument in Lemma~\ref{lem:generic-endpoints}, applied
  to $T_X$ and $T_X^\vee\otimes\omega_X$, supplies sufficiently large
  negative fibre twists; all of them lie in the vector group $\Pic(X)$.
  Apply that lemma with $Y=X$, $B=T_X$, and $\mathcal P=\Pic(X)$.  The
  exponential identification supplies the required Poincar\'e family and all
  negative fibre twists, so the lemma gives simultaneous degree-zero and
  degree-three vanishing on a dense open set.

The jump locus
\[
 J_2=\{M\in\Pic(X):h^2(X,T_X\otimes M)\ge1\}
\]
is a closed analytic subset by upper semicontinuity.  It is proper by (5.4),
so its complement is dense open.  Intersecting the three dense open sets
proves the cohomology assertions in (5.5).  On this intersection,
\[
 \chi(X,T_X\otimes M)=-h^1(X,T_X\otimes M)\le0.
 \tag{5.6}
\]

It remains to evaluate the Euler characteristic.  From (5.3), the rational
groups $H^1(X;\Q)$ and $H^2(X;\Q)$ vanish.  Universal coefficients and
Poincar\'e duality then give $H^4(X;\Q)=0$.  Hence
  $c_1(T_X)=c_1(K_X)=0$ and $c_2(T_X)=0$ in rational cohomology, and
  $c_1(M)=0$ as well.  Since $X$ is
  compact complex, holomorphic Hirzebruch--Riemann--Roch applies \cite{Hir}.  The
degree-six component of
\(\operatorname{ch}(T_X\otimes M)\operatorname{td}(X)\) is consequently
\(\tfrac12c_3(T_X)\), so
\[
 \chi(X,T_X\otimes M)=
 \frac12\left\langle c_3(T_X),[X]\right\rangle.
 \tag{5.7}
\]
Combining (5.6) and (5.7) proves the final inequality.
\end{proof}

